\documentclass[11pt,a4paper]{article}
\usepackage[utf8]{inputenc}
\usepackage[T1]{fontenc}
\usepackage[ngerman]{babel}
\usepackage{helvet}
\renewcommand{\familydefault}{\sfdefault}
\usepackage[left=2cm, right=2cm, top=2cm, bottom=2cm]{geometry}
\usepackage{amsmath,amssymb}
\usepackage{array}
\usepackage{booktabs}
\usepackage{tikz}
\usepackage{pgfplots}
\pgfplotsset{compat=1.18}
\usepackage{fancyhdr}
\usepackage{xcolor}
\usepackage{siunitx}
\sisetup{locale = DE}
\usepackage{lastpage}

% Lösungsfarbe
\definecolor{correct}{RGB}{34,139,94}

\pagestyle{fancy}
\fancyhf{}
\fancyhead[L]{\small Physik Klasse 10E}
\fancyhead[R]{\small\bfseries ERWARTUNGSHORIZONT}
\fancyfoot[C]{\small Seite \thepage\ von \pageref{LastPage}}
\renewcommand{\headrulewidth}{0.4pt}

\setlength{\parindent}{0pt}
\setlength{\fboxsep}{8pt}

% Aufgaben-Befehl
\newcounter{aufgabe}
\newcommand{\aufgabe}[2]{%
\stepcounter{aufgabe}%
\vspace{0.5cm}%
\noindent\fbox{\parbox{\dimexpr\textwidth-2\fboxsep-2\fboxrule}{%
\textbf{Aufgabe \theaufgabe: #1}\hfill\textbf{#2 BE}%
}}%
\vspace{0.3cm}%
}

\newcommand{\lsg}[1]{\textcolor{correct}{#1}}
\newcommand{\be}[1]{\hfill\textcolor{correct}{\small(#1 BE)}}

\begin{document}

% === KOPFBEREICH ===
\noindent\fbox{\parbox{\dimexpr\textwidth-2\fboxsep-2\fboxrule}{
\centering
{\Large\bfseries Klausur: Kinematik}\\[0.2cm]
{\large\bfseries Erwartungshorizont}\\[0.2cm]
{\small Klasse 10E \quad|\quad 90 Minuten \quad|\quad 45 BE}
}}

\vspace{0.4cm}

% === NOTENSPIEGEL (15-NP Sek II MV) ===
\begin{center}
\small
\begin{tabular}{|c|c|c|c|c|c|c|c|c|c|c|c|c|c|c|c|}
\hline
\textbf{NP} & 15 & 14 & 13 & 12 & 11 & 10 & 9 & 8 & 7 & 6 & 5 & 4 & 3 & 2 & 1 \\
\hline
\textbf{ab BE} & 43 & 40,5 & 38 & 36 & 33,5 & 31,5 & 29 & 27 & 24,5 & 22,5 & 20 & 18 & 15 & 12 & 9 \\
\hline
\textbf{ab \%} & 95 & 90 & 85 & 80 & 75 & 70 & 65 & 60 & 55 & 50 & 45 & 40 & 33 & 27 & 20 \\
\hline
\end{tabular}

\vspace{0.2cm}
{\footnotesize NP = Notenpunkte (15-Punkte-Skala, Sek II MV)}
\end{center}

\vspace{0.3cm}

% ═══════════════════════════════════════════════════════════════
% AUFGABE 1: Grundwissen (8 BE, AFB I)
% ═══════════════════════════════════════════════════════════════

\aufgabe{Grundwissen Kinematik}{8}

\begin{enumerate}
\item[a)] \lsg{Die Steigung im $s$-$t$-Diagramm stellt die \textbf{Geschwindigkeit $v$} dar.} \be{1}

\item[b)] \lsg{Die Fläche unter einem $v$-$t$-Graphen stellt den \textbf{zurückgelegten Weg $s$} dar.} \be{1}

\item[c)] Umrechnung: \be{2}

\lsg{$90\,\mathrm{km/h} = 90 \div 3{,}6 = 25\,\mathrm{m/s}$}

\lsg{$54\,\mathrm{km/h} = 54 \div 3{,}6 = 15\,\mathrm{m/s}$}

{\small\color{correct} Je 1 BE pro richtige Umrechnung}

\item[d)] Zuordnung (mit sauberen Linien verbunden): \be{4}

\vspace{0.3cm}

\begin{center}
\begin{tikzpicture}[
    box/.style={draw, minimum width=5cm, minimum height=0.9cm, align=center, font=\small},
    lsgbox/.style={draw, minimum width=4.5cm, minimum height=0.9cm, align=center, font=\small, text=correct}
]
% Linke Kästen (Bewegungsarten) - gemischt wie in Aufgabenstellung
\node[box] (L1) at (0,0) {Gleichförmige Bewegung\\im $s$-$t$-Diagramm};
\node[box] (L2) at (0,-1.4) {Abbremsen\\im $v$-$t$-Diagramm};
\node[box] (L3) at (0,-2.8) {Gleichförmige Bewegung\\im $v$-$t$-Diagramm};
\node[box] (L4) at (0,-4.2) {Beschleunigte Bewegung\\im $s$-$t$-Diagramm};

% Rechte Kästen (Graphenformen) - gemischt wie in Aufgabenstellung
\node[lsgbox] (R1) at (8.5,0) {Waagerechte Linie};
\node[lsgbox] (R2) at (8.5,-1.4) {Nach oben geöffnete Parabel};
\node[lsgbox] (R3) at (8.5,-2.8) {Fallende Gerade};
\node[lsgbox] (R4) at (8.5,-4.2) {Ansteigende Gerade};

% Verbindungslinien (Lösungen)
\draw[correct, thick] (L1.east) -- (R4.west); % Gleichf. s-t → Ansteigende Gerade
\draw[correct, thick] (L2.east) -- (R3.west); % Abbremsen v-t → Fallende Gerade
\draw[correct, thick] (L3.east) -- (R1.west); % Gleichf. v-t → Waagerechte Linie
\draw[correct, thick] (L4.east) -- (R2.west); % Beschl. s-t → Parabel
\end{tikzpicture}
\end{center}

\vspace{0.2cm}
{\small\color{correct} Je 1 BE pro richtige Zuordnung. Linien müssen sauber gezogen sein.}
\end{enumerate}

% ═══════════════════════════════════════════════════════════════
% AUFGABE 2: Freier Fall (8 BE)
% ═══════════════════════════════════════════════════════════════

\aufgabe{Freier Fall -- Reaktionszeit}{6}

\begin{enumerate}
\item[a)] Herleitung der Formel: \be{2}

\lsg{Ausgangsformel: $h = \dfrac{1}{2}g \cdot t^2$}\\[0.2cm]
\lsg{Umstellen nach $t^2$: $t^2 = \dfrac{2h}{g}$} \hfill(1 BE)\\[0.2cm]
\lsg{Wurzel ziehen: $t = \sqrt{\dfrac{2h}{g}}$} \hfill(1 BE)

\item[b)] Berechnung der Reaktionszeit: \be{2}

\lsg{$t = \sqrt{\dfrac{2h}{g}} = \sqrt{\dfrac{2 \cdot 0{,}45\,\mathrm{m}}{10\,\mathrm{m\,s^{-2}}}} = \sqrt{\dfrac{0{,}9\,\mathrm{m}}{10\,\mathrm{m\,s^{-2}}}} = \sqrt{0{,}09\,\mathrm{s^2}}$} \hfill(1 BE)\\[0.2cm]
\lsg{$t = 0{,}3\,\mathrm{s} = 300\,\mathrm{ms}$} \hfill(1 BE)

\item[c)] Erklärung Mond-Experiment: \be{2}

\lsg{Auf dem Mond ist $g$ kleiner ($1{,}6$ statt $10\,\mathrm{m/s^2}$). Das Lineal fällt langsamer und legt bei gleicher Reaktionszeit weniger Strecke zurück (Beispiel: In $0{,}25\,\mathrm{s}$ fällt es auf der Erde $31\,\mathrm{cm}$, auf dem Mond nur $5\,\mathrm{cm}$).} \hfill(1 BE)\\[0.2cm]
\lsg{Die Auflösung des Experiments sinkt: Auf der Erde entsprechen $10\,\mathrm{ms}$ Unterschied ca.\ $2{,}5\,\mathrm{cm}$ auf dem Lineal, auf dem Mond nur ca.\ $0{,}4\,\mathrm{cm}$. Die Ablesegenauigkeit wird schlechter.} \hfill(1 BE)

\vspace{0.3cm}
{\small\color{correct} \textbf{Beispielantwort (volle Punktzahl):} "`Die Fallbeschleunigung auf dem Mond beträgt nur $1{,}6\,\mathrm{m/s^2}$. Das Lineal fällt also viel langsamer und legt für die gleiche Reaktionszeit weniger Strecke zurück. Die Auflösung ist geringer, weil $10\,\mathrm{ms}$ Unterschied nur noch $0{,}4\,\mathrm{cm}$ statt $2{,}5\,\mathrm{cm}$ ausmachen. Das Ablesen wird ungenauer."'}
\end{enumerate}

\newpage

% ═══════════════════════════════════════════════════════════════
% AUFGABE 3: Bremsvorgang (8 BE)
% ═══════════════════════════════════════════════════════════════

\aufgabe{Bremsvorgang}{8}

\begin{enumerate}
\item[a)] \lsg{$v_0 = 108\,\mathrm{km\,h^{-1}} = 108 \div 3{,}6 = 30\,\mathrm{m\,s^{-1}}$} \be{1}

\item[b)] Bremszeit: \be{2}

\lsg{Formel: $v = v_0 + a \cdot t$ mit $v = 0$ $\Rightarrow$ $t = \dfrac{v_0}{|a|}$} \hfill(1 BE)\\[0.2cm]
\lsg{Einsetzen: $t = \dfrac{30\,\mathrm{m\,s^{-1}}}{6\,\mathrm{m\,s^{-2}}} = 5\,\mathrm{s}$} \hfill(1 BE)

\item[c)] Bremsweg: \be{2}

\lsg{Formel: $s = \dfrac{v_0^2}{2|a|}$} \hfill(1 BE)\\[0.2cm]
\lsg{Einsetzen: $s = \dfrac{(30\,\mathrm{m\,s^{-1}})^2}{2 \cdot 6\,\mathrm{m\,s^{-2}}} = \dfrac{900\,\mathrm{m^2\,s^{-2}}}{12\,\mathrm{m\,s^{-2}}} = 75\,\mathrm{m}$} \hfill(1 BE)

\item[d)] Beurteilung: \be{3}

\lsg{\textbf{Ja}, der Fahrer kann rechtzeitig stoppen.} \hfill(1 BE)\\[0.2cm]
\lsg{Begründung: Der Bremsweg beträgt 75 m, das Hindernis ist 80 m entfernt.} \hfill(1 BE)\\[0.2cm]
\lsg{$80\,\mathrm{m} - 75\,\mathrm{m} = 5\,\mathrm{m}$ Sicherheitsabstand verbleiben.} \hfill(1 BE)
\end{enumerate}

% ═══════════════════════════════════════════════════════════════
% AUFGABE 4: Busfahrt (12 BE)
% ═══════════════════════════════════════════════════════════════

\aufgabe{Busfahrt durch Rostock}{12}

\begin{enumerate}
\item[a)] Beschreibung der sechs Phasen: \be{3}

\lsg{Phase 1 (0--4\,s): Der Bus beschleunigt gleichmäßig von 0 auf 12\,m/s.}\\
\lsg{Phase 2 (4--10\,s): Der Bus fährt gleichförmig mit konstant 12\,m/s.}\\
\lsg{Phase 3 (10--14\,s): Der Bus bremst gleichmäßig von 12\,m/s auf 0.}\\
\lsg{Phase 4 (14--18\,s): Der Bus steht (Halt an Ampel/Haltestelle).}\\
\lsg{Phase 5 (18--22\,s): Der Bus beschleunigt wieder von 0 auf 8\,m/s.}\\
\lsg{Phase 6 (22--28\,s): Der Bus fährt gleichförmig mit 8\,m/s.}

{\small\color{correct} Je 0,5 BE pro richtig beschriebene Phase; volle Punktzahl bei 6 korrekten Phasen}

\item[b)] Beschleunigung in Phase 1: \be{2}

\lsg{Formel: $a = \dfrac{\Delta v}{\Delta t}$} \hfill(1 BE)\\[0.2cm]
\lsg{Einsetzen: $a = \dfrac{12\,\mathrm{m\,s^{-1}} - 0}{4\,\mathrm{s} - 0} = \dfrac{12\,\mathrm{m\,s^{-1}}}{4\,\mathrm{s}} = 3\,\mathrm{m\,s^{-2}}$} \hfill(1 BE)

\item[c)] Tabelle der Wegberechnung: \be{3}

\begin{center}
\begin{tabular}{|c|c|l|c|c|}
\hline
\textbf{Phase} & \textbf{Zeit} & \textbf{Wegberechnung (Flächenansatz)} & $\boldsymbol{\Delta s}$ & $\boldsymbol{s_{\textbf{ges}}}$ \\
\hline
1 & 0--4\,s & \lsg{$\Delta s = \frac{1}{2} \cdot 4\,\mathrm{s} \cdot 12\,\mathrm{m\,s^{-1}} = 24\,\mathrm{m}$} & \lsg{24\,m} & \lsg{24\,m} \\
\hline
2 & 4--10\,s & \lsg{$\Delta s = 6\,\mathrm{s} \cdot 12\,\mathrm{m\,s^{-1}} = 72\,\mathrm{m}$} & \lsg{72\,m} & \lsg{96\,m} \\
\hline
3 & 10--14\,s & \lsg{$\Delta s = \frac{1}{2} \cdot 4\,\mathrm{s} \cdot 12\,\mathrm{m\,s^{-1}} = 24\,\mathrm{m}$} & \lsg{24\,m} & \lsg{120\,m} \\
\hline
4 & 14--18\,s & \lsg{$\Delta s = 0$ (Stillstand, $v = 0$)} & \lsg{0\,m} & \lsg{120\,m} \\
\hline
5 & 18--22\,s & \lsg{$\Delta s = \frac{1}{2} \cdot 4\,\mathrm{s} \cdot 8\,\mathrm{m\,s^{-1}} = 16\,\mathrm{m}$} & \lsg{16\,m} & \lsg{136\,m} \\
\hline
6 & 22--28\,s & \lsg{$\Delta s = 6\,\mathrm{s} \cdot 8\,\mathrm{m\,s^{-1}} = 48\,\mathrm{m}$} & \lsg{48\,m} & \lsg{184\,m} \\
\hline
\end{tabular}
\end{center}

\vspace{0.2cm}
{\small\color{correct} \textbf{Bewertung:} 0,5 BE für korrekte Wegberechnung + 0,5 BE für korrektes $\Delta s$ und $s_{\text{ges}}$ pro Zeile (= 6 × 0,5 = 3 BE)}

\newpage

\item[d)] \textbf{Zeichne} das zugehörige $s$-$t$-Diagramm. Beschrifte die $s$-Achse und achte auf die korrekte Kurvenform. \be{4}

\begin{center}
\begin{tikzpicture}
\begin{axis}[
    width=13cm, height=7cm,
    xlabel={$t$ in s},
    ylabel={$s$ in m},
    xmin=0, xmax=30,
    ymin=0, ymax=200,
    grid=both,
    grid style={gray!30},
    major grid style={gray!50},
    axis lines=left,
    xtick={0,4,10,14,18,22,28},
    ytick={0,24,96,120,136,184},
    every axis label/.style={font=\small},
]
% Phase 1: Parabel (beschleunigt) - s = 1.5*t^2
\addplot[very thick, correct, samples=50, domain=0:4] {1.5*x^2};
% Phase 2: Gerade (gleichförmig) - von (4,24) bis (10,96)
\addplot[very thick, correct] coordinates {(4,24) (10,96)};
% Phase 3: Parabel (bremst) - nach unten konkav
\addplot[very thick, correct, samples=50, domain=10:14] {120 - 1.5*(14-x)^2};
% Phase 4: Waagerecht (Stillstand)
\addplot[very thick, correct] coordinates {(14,120) (18,120)};
% Phase 5: Parabel (beschleunigt) - s = 120 + 1*t^2
\addplot[very thick, correct, samples=50, domain=18:22] {120 + 1*(x-18)^2};
% Phase 6: Gerade (gleichförmig)
\addplot[very thick, correct] coordinates {(22,136) (28,184)};

% Markierungspunkte
\addplot[only marks, mark=*, correct, mark size=2pt] coordinates {(0,0) (4,24) (10,96) (14,120) (18,120) (22,136) (28,184)};
% Gestrichelte vertikale Phasengrenzen
\draw[gray!70, dashed, thick] (axis cs:4,0) -- (axis cs:4,200);
\draw[gray!70, dashed, thick] (axis cs:10,0) -- (axis cs:10,200);
\draw[gray!70, dashed, thick] (axis cs:14,0) -- (axis cs:14,200);
\draw[gray!70, dashed, thick] (axis cs:18,0) -- (axis cs:18,200);
\draw[gray!70, dashed, thick] (axis cs:22,0) -- (axis cs:22,200);
\end{axis}
\end{tikzpicture}
\end{center}

\textbf{Bewertungskriterien:}\\
\lsg{Achsenbeschriftung $s$-Achse korrekt ($s$ in m, äquidistante Skalierung, $\geq 184\,\mathrm{m}$)} \hfill(1 BE)\\
\lsg{Parabeln in Phase 1, 3, 5 korrekt (Form und Werte)} \hfill(1,5 BE)\\
\lsg{Geraden in Phase 2 und 6 + waagerechte Linie in Phase 4 korrekt (Werte)} \hfill(1,5 BE)\\[0.2cm]
{\small\color{correct} Sauberkeit und Stetigkeit (knickfreie Übergänge) werden erwartet. Bei Mängeln ggf.\ Abzug.}

\end{enumerate}

\newpage

% ═══════════════════════════════════════════════════════════════
% AUFGABE 5: Überholvorgang (8 BE)
% ═══════════════════════════════════════════════════════════════

\aufgabe{Überholvorgang}{8}

\begin{enumerate}
\item[a)] Umrechnung: \be{1}

\lsg{$v_{\text{Traktor}} = 36\,\mathrm{km\,h^{-1}} = 36 \div 3{,}6 = 10\,\mathrm{m\,s^{-1}}$}\\
\lsg{$v_{\text{PKW}} = 72\,\mathrm{km\,h^{-1}} = 72 \div 3{,}6 = 20\,\mathrm{m\,s^{-1}}$}

\item[b)] Relativgeschwindigkeit: \be{2}

\lsg{Formel: $v_{\text{rel}} = v_{\text{PKW}} - v_{\text{Traktor}}$} \hfill(1 BE)\\[0.2cm]
\lsg{Einsetzen: $v_{\text{rel}} = 20\,\mathrm{m\,s^{-1}} - 10\,\mathrm{m\,s^{-1}} = 10\,\mathrm{m\,s^{-1}}$} \hfill(1 BE)

\item[c)] Zeit bis zum Erreichen: \be{2}

\lsg{Formel: $t = \dfrac{\Delta s}{v_{\text{rel}}}$} \hfill(1 BE)\\[0.2cm]
\lsg{Einsetzen: $t = \dfrac{200\,\mathrm{m}}{10\,\mathrm{m\,s^{-1}}} = 20\,\mathrm{s}$} \hfill(1 BE)

\item[d)] Analyse des Überholvorgangs: \be{3}

\lsg{In $t = 20\,\mathrm{s}$ legt der PKW zurück:}\\
\lsg{$s_{\text{PKW}} = v_{\text{PKW}} \cdot t = 20\,\mathrm{m\,s^{-1}} \cdot 20\,\mathrm{s} = 400\,\mathrm{m}$}\\[0.2cm]
\lsg{Der PKW hat den Traktor nach 20\,s \textbf{erreicht}, aber noch nicht überholt.} \hfill(1 BE)\\[0.2cm]
\lsg{Für sichere Überholung (+ 50\,m Sicherheitsabstand) braucht er zusätzlich:}\\
\lsg{$t_{\text{extra}} = \dfrac{50\,\mathrm{m}}{10\,\mathrm{m\,s^{-1}}} = 5\,\mathrm{s}$}\\
\lsg{Gesamt: $t_{\text{ges}} = 20\,\mathrm{s} + 5\,\mathrm{s} = 25\,\mathrm{s}$}\\[0.2cm]
\lsg{Weg des PKW in 25\,s: $s = 20\,\mathrm{m\,s^{-1}} \cdot 25\,\mathrm{s} = 500\,\mathrm{m}$} \hfill(1 BE)\\[0.2cm]
\lsg{\textbf{Beurteilung:} Der PKW erreicht die Kurve genau nach 500\,m. Das Überholen ist \textbf{kritisch} -- es ist nicht sicher zu empfehlen, da kein Puffer verbleibt.} \hfill(1 BE)

{\small\color{correct} Akzeptiert auch: "`Nicht sicher"', "`Grenzwertig"', "`Zu riskant"'}
\end{enumerate}

% ═══════════════════════════════════════════════════════════════
% AUFGABE 6: Fallschirmspringen (3 BE) - phänomenologisch
% ═══════════════════════════════════════════════════════════════

\aufgabe{Fallschirmspringen}{3}

\begin{enumerate}
\item[a)] Warum wird der Springer immer schneller? \be{1}

\lsg{Akzeptierte Antworten (eine davon genügt):}
\begin{itemize}
\item \lsg{Die Erde zieht ihn nach unten / Die Schwerkraft zieht ihn an}
\item \lsg{Die Erdanziehung / Gravitation beschleunigt ihn}
\item \lsg{Die Gewichtskraft wirkt nach unten}
\item \lsg{Er fällt frei, weil nichts ihn hält/bremst}
\end{itemize}

\item[b)] Konstante Fallgeschwindigkeit ohne Schirm: \be{1}

\lsg{Akzeptierte Antworten (eine davon genügt):}
\begin{itemize}
\item \lsg{Der Luftwiderstand wird so groß wie die Schwerkraft}
\item \lsg{Die Kräfte heben sich auf / sind gleich groß (Gleichgewicht)}
\item \lsg{Der Luftwiderstand bremst ihn so stark, dass er nicht schneller wird}
\item \lsg{Die Luft bremst ihn immer stärker, bis er nicht mehr beschleunigt}
\end{itemize}

\item[c)] Geringere Grenzgeschwindigkeit mit Schirm: \be{1}

\lsg{Akzeptierte Antworten (eine davon genügt):}
\begin{itemize}
\item \lsg{Der Schirm ist größer → mehr Luftwiderstand}
\item \lsg{Größere Fläche = mehr Bremswirkung}
\item \lsg{Der Schirm fängt mehr Luft}
\item \lsg{Das Gleichgewicht wird schon bei niedriger Geschwindigkeit erreicht}
\end{itemize}
\end{enumerate}

\vspace{0.5cm}
\hrule
\vspace{0.5cm}

% === AFB-ZUSAMMENFASSUNG ===

\begin{center}
\fbox{\parbox{0.9\textwidth}{
\centering
\textbf{Zusammenfassung der Anforderungsbereiche}\\[0.3cm]
\begin{tabular}{|l|c|c|c|}
\hline
\textbf{AFB} & \textbf{Teilaufgaben} & \textbf{BE} & \textbf{Anteil} \\
\hline
I (Reproduktion) & 1a-d, 2a, 3a, 5a & 12 & 27\% \\
\hline
II (Anwendung) & 2b, 3b-c, 4a-c, 5b-c, 6a & 20 & 44\% \\
\hline
III (Transfer) & 2c, 3d, 4d, 5d, 6b-c & 13 & 29\% \\
\hline
\textbf{Gesamt} & & \textbf{45} & 100\% \\
\hline
\end{tabular}
}}
\end{center}

\vspace{0.3cm}

\begin{center}
\textit{--- Ende des Erwartungshorizonts ---}
\end{center}

\end{document}
